\SourceLeaf{62}
Pars prima negotii propositi confecta est, circumscriptis formula finita ipsius $(1+n\cos\varphi)^{-1}$ coefficientibus. Restat altera, ut videamus, qua ratione a potestatis alicuius coefficientibus datis transire liceat ad quaesitos potestatis proxime inferioris.

$\delta$) Facillime hoc negotium ita perficitur, ut in expressionibus, quibus tum coefficiens seriei indefinitae pro $(1+n\cos\varphi)^{\nu}$ initialis, tum generalis exhibentur, in locum ipsius $\nu$ substituatur $\nu-1$, et tunc eliciatur resultantium cum anterioribus relatio.

I. Erat (6) ipsius coefficientis initialis $A$, in seriem evoluti, terminus initialis $1$, qui igitur substitutione non adficitur. Erat generalis $h^{\text{tus}}$
\[
= \frac{1}{2^{2h}}\,{}^{\nu}\mathfrak{B}^{(2h)} .  {}^{2h}\mathfrak{B}^{(h)}\,n^{2h},
\]
qui, substituendo $\nu-1$ loco $\nu$, mutatur in:
\[
\frac{1}{2^{2h}}\,{}^{\nu-1}\mathfrak{B}^{(2h)} .  {}^{2h}\mathfrak{B}^{(h)}\,n^{2h}.
\]
Manifesto igitur nascetur ex priore posterior, si in
\[
\frac{\nu-2h}{\nu} = 1-\frac{2h}{\nu}
\]
ducatur ille; adeoque, cum idem sit, terminum generalem ipsius $A$ $h^{\text{tum}}$ in $\frac{2h}{\nu}$ ducere, ac eundem, secundum $n$ differentiatum, per $\nu\,\mathrm{d}n$ dividere et per $n$ multiplicare, habebimus, si vocetur $B$ ipsius $(1+n\cos\varphi)^{\nu-1}$ terminus initialis, generatim
\[
B = A - \frac{n\,{}^{*}\!\mathrm{d}A}{\nu\,\mathrm{d}n}.
\]
\par\noindent\hspace*{\parindent}II. Erat coefficientis indefiniti ad $(1+n\cos\varphi)^{\nu}$ pertinentis $A^{(r)}$ terminus initialis
\[
\frac{1}{2^{r-1}}\,{}^{\nu}\mathfrak{B}^{(r)}\,n^{r},
\]
fit itaque, si vocetur $B^{(r)}$ qui, ad $(1+n\cos\varphi)^{\nu-1}$ pertinens, elicitur substituendo in istum $\nu-1$ loco $\nu$, huius terminus initialis
\[
\frac{1}{2^{r-1}}\,{}^{\nu-1}\mathfrak{B}^{(r)}\,n^{r}.
\]
Hic vero, cum obtineatur illum per $\frac{\nu-r}{\nu}$ multiplicando sive, quod idem est,
